% ============================================================
% DISCUSSION
% ============================================================

\section{Discussion}
\label{sec:discussion}

The results of our comparative evaluation contrast with positive results reported in prior reinforcement learning portfolio management studies. Mean-Variance optimization achieved 366\% cumulative return over the test period, more than double the next-best strategy, and outperformed all three RL agents on both absolute and risk-adjusted metrics. This section interprets these findings, situates them within the broader literature, acknowledges limitations, and identifies directions for future research.

% ------------------------------------------------------------
\subsection{Interpreting Mean-Variance Dominance}
\label{sec:mvo_analysis}
% ------------------------------------------------------------

The strong performance of Mean-Variance optimization admits several complementary explanations rooted in both market conditions and algorithmic characteristics. The test period spanning January 2024 through October 2025 coincided with a sustained cryptocurrency bull market following the prolonged bear market of 2022--2023. In trending market regimes, momentum signals embedded in rolling return estimates can carry predictive information, as recent winners tend to continue outperforming in the short term \cite{jegadeesh1993momentum}. The 60-day estimation window employed by our MVO baseline appears to have captured these momentum dynamics effectively, enabling the optimizer to concentrate allocations in assets with strong recent performance.

Beyond favorable market conditions, Mean-Variance optimization enjoys substantial sample efficiency advantages over reinforcement learning approaches. Where our RL agents required hundreds to thousands of training episodes to learn allocation policies through trial-and-error interaction, MVO directly estimates the sufficient statistics---expected returns and covariances---from historical data and solves a closed-form optimization problem. This estimation-based approach uses available data efficiently, whereas model-free RL must rediscover these statistical relationships through gradient descent on noisy reward signals. The Ledoit-Wolf shrinkage estimator~\citep{ledoit2004honey} further enhances MVO's robustness by addressing the instability of sample covariance matrices in high-dimensional settings \cite{ledoit2004honey}, regularizing estimates toward a structured target and reducing sensitivity to outliers.

We caution, however, that MVO's dominance may not generalize beyond the specific market conditions of our test period. Classical portfolio theory assumes stationary return distributions, an assumption violated in cryptocurrency markets characterized by regime shifts, bubble dynamics, and structural breaks \cite{james2022correlations}. In mean-reverting or highly volatile regimes where historical estimates lose predictive validity, the comparative advantage of MVO would likely diminish. The strong performance observed here should be interpreted as evidence that simple estimation-based methods remain competitive baselines rather than as a general indictment of learning-based approaches.

% ------------------------------------------------------------
\subsection{Explaining Reinforcement Learning Agent Behavior}
\label{sec:rl_analysis}
% ------------------------------------------------------------

The behavior of our reinforcement learning agents reveals instructive patterns about the challenges of applying deep RL to portfolio optimization. REINFORCE achieved performance nearly identical to the Equal Weight baseline (166.98\% vs.\ 168.50\% return), with portfolio allocations that converged to near-uniform distributions across tradable assets. This convergence suggests that the policy gradient learned to minimize transaction costs rather than exploit predictive signals in the observation space. Given the reward structure penalizing turnover at 0.1\% per unit rebalancing, a risk-averse policy that maintains stable allocations and avoids trading costs represents a local optimum that gradient descent may struggle to escape. This interpretation aligns with the theoretical analysis of~\citet{demiguel2009optimal}, who demonstrated that when estimation error dominates, naive diversification can outperform sophisticated optimization by avoiding costly mistakes.

The value-based agents underperformed all baselines. Both DQN (155.47\%) and DDQN (144.50\%) achieved lower returns than the three baseline strategies despite systematic hyperparameter optimization. Their elevated turnover rates of 10.88\% and 8.39\% respectively suggest that the delta-based action space, while intuitive for human interpretation, may have encouraged excessive rebalancing. Each action in our 70-action catalog applies a discrete portfolio adjustment, and the Q-network must learn which adjustments are profitable in which contexts---a challenging credit assignment problem when rewards are noisy and delayed. The observation that DDQN performed worse than standard DQN, contrary to expectations from the broader RL literature~\citep{vanHasselt2016ddqn}, indicates that reduced Q-value overestimation may have produced overly conservative policies in this domain, where some degree of optimism could encourage beneficial exploration.

Sample complexity poses a significant challenge for value-based methods in our setting. The canonical state representation spans 8,917 dimensions before projection, and the training budget of 400--500 episodes provides limited coverage of the state-action space. Prior work applying DQN to portfolio management~\citep{jiang2016drlt, lucarelli2020dqlcrypto} employed smaller asset universes and shorter lookback windows, reducing dimensionality and potentially easing the learning problem. Scaling deep RL to larger, more realistic portfolio settings may require architectural innovations such as attention mechanisms for cross-asset reasoning or hierarchical policies that separate regime detection from allocation decisions.

% ------------------------------------------------------------
\subsection{Comparison with Prior Work}
\label{sec:comparison_prior}
% ------------------------------------------------------------

Our findings contrast with several influential papers reporting strong RL outperformance in portfolio management. \citet{jiang2017eIIE} demonstrated that their EIIE architecture substantially outperformed baselines on cryptocurrency data from 2015--2017, a period characterized by the initial Bitcoin bull run and the emergence of altcoin markets. \citet{ye2020sarl} achieved similar success with state-augmented RL on Chinese equity markets. However, direct comparison is complicated by significant methodological differences: these studies employed smaller asset universes (12 assets vs.\ our 37), shorter evaluation periods, and in some cases did not impose realistic transaction cost structures.

The divergence between our results and prior claims admits several interpretations. First, market regime matters profoundly---the cryptocurrency landscape of 2024--2025 differs substantially from 2015--2017 in terms of market maturity, institutional participation, and correlation structure. Strategies that exploited inefficiencies in nascent markets may not transfer to more efficient conditions. Second, publication bias likely inflates reported RL performance, as negative results are systematically underrepresented in the literature. Our regime-diverse validation protocol and extended out-of-sample test period impose stricter evaluation standards than typical in this literature. Third, the fully-invested constraint in our formulation prevents agents from avoiding exposure during drawdowns, whereas some prior work permitted cash allocations that could trivially improve risk-adjusted metrics.

% ------------------------------------------------------------
\subsection{Limitations}
\label{sec:limitations}
% ------------------------------------------------------------

Several limitations temper the conclusions that can be drawn from this study. Most significantly, our test period represents a single market regime---a post-bear-market recovery transitioning into a bull market. Results from this favorable environment may not generalize to sustained bear markets, high-volatility crash periods, or prolonged sideways consolidation. A comprehensive evaluation would require testing across multiple complete market cycles, which the relatively short history of liquid cryptocurrency markets does not yet permit.

Our transaction cost model, while more realistic than many prior studies, remains simplified. The fixed 0.1\% cost per unit turnover ignores slippage, market impact, and the spread variation that characterizes cryptocurrency exchanges. Large rebalancing trades in less liquid assets would incur substantially higher implicit costs than our model assumes, potentially favoring concentrated positions in major assets. Additionally, our agents were frozen after training and evaluated without online adaptation. Continual learning approaches that update policies as market conditions evolve could narrow the gap with estimation-based methods that naturally incorporate recent data.

The hyperparameter search, while systematic, operated under computational constraints. Fifty Optuna trials may not have discovered optimal configurations for agents with high-dimensional hyperparameter spaces, and the episode budgets allocated for production training may have been insufficient for convergence, particularly for the value-based methods struggling with sample complexity. Finally, while Mean-Variance optimization substantially outperformed RL agents, it benefits from a form of hindsight through its rolling estimation window---returns observed after portfolio formation but within the estimation window for subsequent days implicitly inform allocations, a subtle advantage not available to RL agents trained on strictly historical data.

% ------------------------------------------------------------
\subsection{Implications and Future Directions}
\label{sec:future}
% ------------------------------------------------------------

Our findings carry implications for both practitioners and researchers working at the intersection of reinforcement learning and quantitative finance. For practitioners, the results underscore the importance of including simple baselines in any evaluation of learning-based portfolio strategies. The strong performance of Mean-Variance optimization---a technique dating to the 1950s---relative to deep RL architectures indicates that more complex methods do not automatically outperform classical approaches in financial domains. Any deployment of RL-based allocation systems should demonstrate clear outperformance against classical alternatives under realistic transaction cost assumptions and across diverse market conditions.

For the research community, our results highlight several promising directions for advancing RL portfolio management. The sample complexity challenges faced by our value-based agents suggest that architectural innovations enabling more efficient learning deserve priority. Transformer architectures with cross-asset attention mechanisms could enable agents to reason about correlation structure without explicit estimation, potentially combining the strengths of MVO and deep learning. Hierarchical reinforcement learning offers another avenue, decomposing the allocation problem into regime detection at a slow timescale and tactical positioning at a fast timescale. Such decomposition could improve credit assignment and enable transfer learning across market regimes. Finally, online and continual learning approaches that adapt policies to evolving market conditions could address the non-stationarity that fundamentally limits offline-trained agents. The gap between research benchmarks and practical deployment remains substantial, and closing it will require sustained attention to evaluation methodology, computational efficiency, and robustness to distribution shift.
